% displayname : Mécanique quantique ondulatoire
\documentclass[12pt,fancy]{cours}
\Chapit{6}{7}
\begin{document}
\maketitre[***]{
  En 1905, Einstein propose l'idée de quantum de lumière. L'onde électromagnétique est alors associée à un corpuscule, le photon. De Broglie propose en 1923 d'étendre cette idée de dualité onde-corpuscule à la matière. C'est le point de départ d'une aventure intellectuelle qui donna naissance à la mécanique quantique. Dans ce chapitre, nous présenterons les fondements de la mécanique quantique ondulatoire développés essentiellement entre 1923 et 1927 (par de Broglie, Schrödinger, Heisenberg, Dirac, Pauli, Hilbert, Born ...). La problème de la mesure en mécanique quantique, par contre, ne fait pas partie du programme.
}{
\item Cours de Physique quantique - 0
\item Document de cours Physique quantique - 1
\item TD Physique quantique
}
%%%%%%%%%%%%%%%%%%%%%%%%%%%%%%%%%%%%%%%%%%%%%%%%%%%%%%%%%%%%%%%%%%%%%%%%%%


\section{Description ondulatoire probabiliste}
\label{sec:description_proba}
Dans tout ce chapitre, afin de simplifier les notations nous allons nous placer à une dimension d'espace. La généralisation des concepts évoqués ci-dessous à des dimensions supérieures ne pose pas de difficulté particulière. Nous étudions ici la théorie quantique des particules \imp{massives} de masse $m$ dans la limite non relativiste $v\ll c$. L'étude du photon, en particulier, ne fait pas partie du programme. 
\subsection{Dualité onde-corpuscule}
\label{subsec:interfs}

\parag{Interférences des ondes de matière} Nous étudions une expérience d'interférence des fentes de Young réalisée en 1992 par Shimizu-Shimizu-Takuma. Un nuage de quelques millions d’atomes de néon est d’abord capturé et refroidi au millikelvin dans
un piège laser. Les atomes sont ensuite lâchés, sans vitesse
initiale, au-dessus d’un écran percé de
deux fentes parallèles, et tombent sous l'effet de la pesanteur. Les fentes sont
séparées d'une distance $a = \SI{6,0}{\micro m}$, et sont placées devant un écran à une distance $D=\SI{85}{cm}$. La plaque détectrice enregistre l’impact de chaque atome, un impact étant représenté par un point sur la figure.

\begin{figure}[h!]
\centering
\includegraphics[width=0.6\textwidth]{shimizu2.png}
\caption{Expérience de Shimizu-Shimizu-Takuma. Interférences d'atomes froids avec une bifente de Young.}
\label{fig:interfs_electrons_schema}
\end{figure}
% On observe alors les résultats suivants.
\begin{enumerate}
  \item  Chaque électron produit un point d'impact bien défini. On en déduit que les électrons interagissent avec l'écran comme le ferait un \imp{corpuscule}, i.e. une particule classique ponctuelle. 
  %, parce que les électrons n'ont pas tous la màme vitesse initiale.
  \item En envoyant un grand nombre d'atomes, on observe une \important{figure d'interférences} identique à celle qu'on aurait obtenu avec de la lumière : la distribution des impacts sur l'écran est celle d'une \imp{onde}.
  \item La position des points d'impact est \important{aléatoire}. Ce caractère aléatoire ne vient pas simplement du manque de connaissance sur le trajet exact des atomes, mais est une propriété intrinsèque des particules quantiques.
\end{enumerate}
% pour chaque particule  La situation devient plus étonnante lorsqu'on laisse l'expérience se dérouler avec un grand nombre d'électrons. Le point d'impact de chaque électron est  sur le capteur, et il n'est pas possible de prédire avec certitude ce point d'impact même en connaissant parfaitement les conditions initiales. Lorsque le nombre d'impacts s'accumule sur le détecteur ), on Chaque impact est aléatoire, mais l'ensemble des impacts se distribue selon une loi identique à la figure d'interférences qu'on observerait avec des ondes lumineuses avec $\lambda=h/p$.
Une description \important{probabiliste et ondulatoire} des particules s'impose: c'est le saut conceptuel franchi par L. de Broglie (1924). On parle de dualité onde-corpuscule.

\begin{prop}[Dualité onde-corpuscule]\label{prop:}
Une particule quantique n'est ni un corpuscule, ni une onde classique, mais peut présenter des comportements à la fois ondulatoire et corpusculaire. Une particule est décrite par une \imp{onde quantique}.
\end{prop}

Le dictionnaire entre les grandeurs corpusculaires et ondulatoires se fait grâce à la constante de Planck réduite $\hbar$. Nous connaissons déjà les expressions ondulatoires des grandeurs mécaniques énergie $E$ et quantité de mouvement  $\Vec{p}$ pour un photon de pulsation $\omega$ et vecteur d'onde $\Vec{k}$. Il s'agit de la relation de Planck-Einstein $E=\hbar\omega$ et de la relation de de Broglie $\Vec{p}=\hbar\Vec{k}$. En vertu de la dualité onde-corpuscule, ces relations s'appliquent aussi aux particules massives.

\begin{theorem}[Relations de dualité]\label{form:}
  À toute particule quantique d'énergie $E$ et de quantité de mouvement $\Vec{p}$ est associée une onde de pulsation $\omega$ et de vecteur d'onde $\Vec{k}$ par les relations
  $$
  \begin{cases}
    E=\hbar\omega\quad (\text{Planck-Einstein}),\\
    \Vec{p}=\hbar\Vec{k} \quad (\text{de Broglie}).
  \end{cases}
  $$
  La constante $\hbar = \frac{h}{2\pi} = \SI{1.05e-34}{J.s}$ est la constante de Planck réduite.
\end{theorem}

\begin{exemple}{Longueur d'onde de de Broglie}
  Considérons une particule libre de masse $m$ et décrite classiquement par sa vitesse $v$. On donne la constante de Planck $h = \SI{6.63e-34}{J.s}$.
  \begin{quest}
    \item	Exprimer la quantité de mouvement $\Vec{p}$ en fonction de la masse $m$ et de la $v$. En déduire l'expression de la longueur d'onde de de Broglie $\lambda$ de l'onde quantique associée à cette particule. 
    \item Exprimer l'énergie $E$ de la particule en fonction de sa masse $m$ et de sa vitesse $v$. En déduire l'expression de la fréquence de Planck-Einstein $\nu$ de l'onde quantique associée à cette particule.
    \item \AN~:électron de masse $m = \SI{9.1e-31}{kg}$ et de  vitesse $v \simeq \SI{e6}{m.s^{-1}}$, ballon de football de masse $m= \SI{0.5}{kg}$ et de vitesse $v \simeq \SI{30}{m.s^{-1}}$. Commenter.
  \end{quest}
  
  \examplesol{
  % Considérons une particule de vitesse $\Vec{v}$, d’impulsion $\Vec{p}$ et d’énergie $E$. À cette particule est associée une onde de vecteur d’onde $\Vec{k}$ et de pulsation $\omega$, suivant le \Tab{dico_Ep}.
  \begin{quest}
    \item L’impulsion d’une particule est donnée par la relation de \imp{De Broglie}
    \begin{equation}
      \Vec{p} = \hbar \Vec{k}.
    \end{equation}
    Or un corpuscule vérifie $\Vec{p} = m \Vec{v}$ et donc la longueur d’onde $\lambda$ vérifie
    \begin{equation}
       mv = \frac{h}{2\pi}\times \frac{2\pi}{\lambda} = \frac{h}{\lambda} \Rar \boxed{\lambda = \frac{h}{p} = \frac{h}{mv}}.
    \end{equation}
    \item  L’énergie d’une particule libre vérifie la relation de \imp{Planck-Einstein}
    \begin{equation}
      E=\hbar \omega.
    \end{equation}
    Or un corpuscule libre vérifie $E=\frac{1}{2}mv^2$, et donc la fréquence $\nu$ vérifie
    \begin{equation}
      \frac{1}{2}mv^2 = \frac{h}{2\pi}\times 2\pi \nu = h \nu \Rar \boxed{\nu = \frac{mv^2}{2h}}.
    \end{equation}
  \end{quest}
  }
\end{exemple}

\begin{exemple}{Interfrange dans l’expérience de Shizimu-Shimizu-Takuma}
  Pour un atome froid d'He de masse $m = \SI{6.64e-27}{kg}$ de vitesse $\SI{2.4}{m.s^{-1}}$, déterminer la longueur d'onde de de Broglie $\lambda$. En déduire l'interfrange $i$ de la figure d'interférences mesuré dans l'expérience de Shimizu-Shimizu-Takuma.

\end{exemple}
\examplesol{
  $\lambda = \SI{8.5}{nm}$ et  $i = \frac{\lambda D}{a} = 8.5e-9/6e-6*85e-2$.
}

% \begin{exemple}{Expérience de Davisson-Germer}
%   En 1927, les physiciens américains Davisson et Germer fournissaient la preuve expérimentale de l’hypothèse de Louis De
% Broglie en mettant en évidence le phénomène de diffraction d’électrons sur un échantillon monocristallin de nickel. Quelques
% mois plus tard, le britannique G.P Thomson confirmait ce résultat en faisant passer un faisceau d’électrons monocinétique à
% travers une mince feuille de métal. Avec des électrons accélérés par une
% différence de potentiel (tension) de l’ordre du kilovolt (kV), il obtint sur une
% plaque photographique placée derrière la cible une figure de diffraction
% identique à celle observée avec des rayons X de même énergie.
% \begin{quest}
%   \item	En quoi l’expérience de G.P Thomson confirma-t-elle la nature ondulatoire des électrons ?
%   \item Donner l’ordre de grandeur de la longueur d’onde des rayons X. L’utilisation de ces derniers vous semble-t-elle adaptée pour mener une étude cristallographique de diffraction ? Justifier.
%   \item Montrer que la longueur d'onde de de Broglie des électrons utilisés par Thomson vaut $\lambda =\frac{\num{1.23}}{\sqrt{U}}$ nm où $U$ est la tension accélératrice en V. En déduire qu'elle correspond bien à la longueur d'onde des rayons X.
% \end{quest}
% \examplesol{}
% \end{exemple}

\begin{table}[h!]
  \centering
\renewcommand{\arraystretch}{1.5}
  \begin{tabular}{|c|c|c|}
    \hline
Grandeur mécanique &   Corpuscule (classique) & Onde (quantique) \\
    \hline
  Quantité de mouvement  & $\Vec{p}=m \Vec{v}$ & $\Vec{p} = \hbar \Vec{k}$\\
    \hline
    % \vspace{0.25cm}
  Énergie cinétique &  $E=\frac{1}{2} m v^2$ & $E=\hbar \omega$
  \\%[0.25cm]
    \hline
  \end{tabular}
  \caption{Dictionnaire entre les expressions corpusculaires, classiques (en fonction de la vitesse $\Vec{v}$) et les expressions ondulatoires, quantiques (en fonction du vecteur d'onde $\Vec{k}$ et de la pulsation $\omega$) des grandeurs mécaniques impulsion $\Vec{p}$ et énergie $E$, pour une particule quantique massive de masse $m$.}
  \label{tab:dico_Ep}

\end{table}


% \begin{figure}[h!]
% \centering
% \includegraphics[width=0.8\textwidth]{double_slit_wiki.pdf}
% \caption{Évolution de la figure d'interférences au cours du temps [d'après l'article Wikipédia \guill{fentes de Young}]. Le cercle rouge sur la figure a) indique un impact d'électron sur le détecteur. Chaque impact est aléatoire mais l'ensemble se répartit pour former une figure d'interférences avec des franges bien visibles en d). L'interfrange mesuré vaut $i=\frac{\lambda D}{a}$ où $D$ est la distance entre l'écran (le détecteur) et les fentes, $a$ la distance entre les bifentes et $\lambda=h/p$ la longueur d'onde de de Broglie des ondes de matière.}
% \label{fig:bifentes_e-}
% \end{figure}

\parag{Un phénomène probabiliste non classique}
Le caractère aléatoire de la position du point d’impact des atomes ne provient pas d’un manque d’information sur les conditions initiales de l’électron : la mécanique quantique est une théorie \imp{intrinsèquement probabiliste}. La distribution des points d’impact sur l’écran est donnée par une certaine \imp{densité de probabilité} $\rho(x)$, telle que $\d P = \rho(x) \d x$ représente la probabilité que l’électron impacte l’écran à la position $x$.


Mais ce n’est pas une théorie probabiliste au sens classique du terme, comme le montre l’expérience d’interférences. En effet, Si l'on bouche une fente, on obtient une certaine densité de probabilité $\rho_1(x)$. Si on bouche l’autre fente, on obtient la distribution $\rho_2(x)$ similaire à $\rho_1(x)$, mais translatée de la distance $a$. En supposant, comme en physique classique, que l’électron \textit{doit} passer ou bien par la première fente ou bien par la deuxième, et que ces deux événements sont indépendants, la densité de probabilité avec les deux fentes ouvertes serait
\begin{equation}
  \rho_{\rm classique}(x) = \rho_1(x) + \rho_2(x).
\end{equation}

Mais lorsqu'on réalise l'expérience, on \textit{constate} que la densité de probabilité obtenue lorsque les deux fentes sont ouvertes fait apparaître un terme d’interférences,
\begin{equation}
  \rho(x) = \rho_1(x) + \rho_2(x) + \rho_{12}(x) \neq \rho_1(x) + \rho_2(x).
\end{equation}
 Autrement dit, lorsqu'on observe les interférences, cela n'a \textit{pas de sens} de vouloir savoir par quelle fente est passée la particule. Et cela est logique si l’on raisonne en terme ondulatoire : lorsqu’on observe des interférences lumineuses dans le dispositif des fentes de Young, on interprète l’expérience par une division du front d’onde de la lumière. Une chose similaire se passe ici.

 \begin{figure}[H]
  \centering
  \subfigure[]{\label{fig:}
    
  \includegraphics[width=0.4\textwidth]{interf_classique.pdf}
  }\hfill
  \subfigure[]{\label{fig:}
    
  \includegraphics[width=0.4\textwidth]{interf_quantique.pdf}
  }
  \caption{
    (a) Distribution de probabilité classique $\rho_{\rm classique}(x) = \rho_1(x) + \rho_2(x)$ obtenue en supposant que l'électron passe par une seule fente à la fois.
    (b) Distribution de probabilité quantique $\rho(x)$ obtenue. On observe un terme d'interférences $\rho_{12}(x)$ tel que $\rho(x) = \rho_1(x) + \rho_2(x) + \rho_{12}(x)$.
  }
  \label{fig:}
 \end{figure}


% \begin{remarque}

%  Conceptuellement c'est assez déroutant si l'on réfléchit en terme de corpuscule: en ouvrant une fente supplémentaire, on a empêché les atomes d'atteindre certaines régions\footnote{Comme le dit J.L.Basdevant: \guill{C'est incroyable de pouvoir empêcher les gens d'entrer chez vous en ouvrant une deuxième porte}.}. En revanche la description mathématique en terme d'onde est très claire (voir Sec.~\ref{subsec:fonction_onde}). On rencontre parfois l’interprétation que l’électron passe par les deux fentes à la fois, ce qui est très maladroit comme formulation : à éviter, donc !
% \end{remarque}

\subsection{La fonction d'onde}

\parag{Définition}
Comparons aux expériences d’interférences lumineuses.
Dans le modèle scalaire de la lumière, on décrit la lumière par une vibration lumineuse $s(x,t)$, qui est un champ scalaire réel, fonction de la position $x$ et du temps $t$.
Par analogie avec l’optique, la mécanique quantique postule qu’une particule est décrite par une \imp{fonction d’onde}.

\begin{definition}[Fonction d’onde]\label{def:}
  Une particule quantique est décrite par un \imp{champ scalaire complexe} $\psi(x,t)$ : la fonction d’onde.
\end{definition}

Néanmoins ce champ n'est pas à valeurs réelles mais à valeurs \important{complexes}. Habituellement les complexes ne sont qu'un outil de calcul et seule la partie réelle nous intéresse. La situation est très différente en mécanique quantique où le caractère complexe de la fonction d'onde est \textit{fondamental}. Une fonction d’onde possède donc un module $|\psi|(x,t)$ et une phase $\vphi(x,t)$, soit
\begin{equation}
  \psi = |\psi|\,e^{i\vphi}.
\end{equation}

\parag{Propriétés de la fonction d’onde}
%
% De plus, l’intensité lumineuse est définie par la moyenne temporelle quadratique
% \begin{equation}
%   I = \moy{s^2}.
% \end{equation}
% Lorsque les conditions d’interférences sont réunies, l’intensité résultant de la superposition de deux ondes lumineuses vérifie
% \begin{equation}
%   I = \avg{s_1+s_2}^{2}=I_1+I_2+2\sqrt{I_1I_2}\cos(\Delta\varphi)
% \end{equation}

% \label{subsec:fonction_onde}
%Une façon de se convaincre de la nécessité du caractère complexe de la fonction d'onde est de raisonner par l'absurde avec une fonction d'onde réelle harmonique pour une particule libre de la forme $\psi(x)=A\cos(\frac{2\pi x}{\lambda})=A\cos(\frac{px}{\hbar})$. Le problème est qu'un tel choix viole le principe d'incertitude de Heisenberg qui indique que si l'impulsion $p$ est exactement connue, on ne devrait avoir aucune information à propos de la position de la particule.
% }.
%
% Il s'agit du premier postulat de la mécanique quantique: la particule est décrite par un champ scalaire complexe noté $\psi$ appelé \important{fonction d'onde}, tel que
% \begin{equation}
% (x,t)\in\mathbb{R}\times\mathbb{R}\mapsto \psi(x,t)\in\mathbb{C}
% \end{equation}
 La fonction d'onde $\psi(x,t)$ représente l'\important{amplitude de probabilité} (complexe) de trouver la particule à l'abscisse $x$ à l'instant $t$. Cela signifie que $|\psi(x,t)|^2$ représente une densité de probabilité.
 \begin{prop}[Interprétation probabiliste de la fonction d’onde]\label{prop:}

   La \imp{densité de probabilité} de présence (en $\SI{}{m^{-1}}$) de la particule à l’abscisse $x$ à l'instant $t$ vaut
   \begin{equation*}
     \rho(x,t)=\abs{\psi(x,t)}^2.
   \end{equation*}
   La grandeur $\d P(x,t)=\rho(x,t)\d x$
donne la probabilité de présence de la particule dans un intervalle de longueur $\d x$ autour de l’abscisse $x$ à l’instant $t$.
 \end{prop}
La fonction d'onde a une dimension physique qui dépend de la dimension d'espace étudiée. À une dimension, $[\rho] = L^{-1}$  et donc $[\psi] = L^{-1/2}$ : la fonction d’onde est homogène à des $m^{-1/2}$.
La probabilité de trouver la particule dans tout l'espace vaut $1$, par définition d'une probabilité. Cela impose une condition de normalisation de la fonction d'onde. On a en effet 
\begin{equation}
  P_{\rm espace} = 1 = \int_{-\infty}^{+\infty} \d P = \int_{-\infty}^{+\infty}\abs{\psi(x,t)}^2\,\d x.
\end{equation}

\begin{prop}[Normalisation de la fonction d’onde]\label{prop:}
  La probabilité totale de trouver la particule dans tout l’espace vaut
  \begin{equation*}
    \int_{-\infty}^{+\infty}\abs{\psi(x,t)}^2\,\d x=1.
  \end{equation*}
\end{prop}

\begin{exemple}{Retour sur l'expérience d'interférences}
Nous verrons en Sec.~\ref{subsec:schrodinger} que l'équation d'évolution de la fonction d'onde, l'équation de Schrödinger, est une équation linéaire. Ainsi la fonction d’onde obéit au \imp{principe de superposition}, comme en optique. On peut alors décrire simplement les résultats de l'expérience d'interférences des ondes de matière. Notons $\psi_1(x,t)$ la fonction d’onde d’un électron issue de la fente $1$. Idem pour $\psi_2(x,t)$ et la fente $2$. La fonction d’onde résultante sur l’écran vaut $\psi(x,t) = \psi_1(x,t) + \psi_2(x,t)$. 
\begin{quest}
  \item	
  Calculer la densité de probabilité de présence de la particule $\rho(x,t)$ sur l'écran et identifier le terme d'interférences. 
  \item Interpréter l'expérience de Shimizu-Shimizu-Takuma.
\end{quest}

\examplesol{
  La densité de probabilité de détection d'une particule en $x$ à l'instant $t$ est donnée par
$$
\rho^2=\abs{\psi_1+\psi_2}^2=\abs{\psi_1}^2+\abs{\psi_2}^2+\left(\psi_1\psi_2^{\star}+\psi_1^{\star}\psi_2\right) = \rho_1+\rho_2 + \rho_{12}
%}_{\text{terme d'interférences}}
$$
où le terme d’interférences vaut \fbox{$\rho_{12}=\psi_1\psi_2^{\star}+\psi_1^{\star}\psi_2$}. On retrouve une formule d’interférences analogue à celle rencontrée en optique pour l’intensité lumineuse $I$.
}
\end{exemple}
% \begin{remarque}
% Ce résultat est à rapprocher de la formule de Fresnel en optique où $\avg{A_1(x,t)+A_2(x,t)}^{2}=I_1(x)+I_2(x)+2\sqrt{I_1(x)I_2(x)}\cos(\Delta\varphi(x))$ qui diffère de $I_1(x)+I_2(x)$.
% \end{remarque}

% \subsection{Importance de la statistique}




\subsection{Relation d'indétermination de Heisenberg}
\label{subsec:heisenberg}

Une particule quantique n’est donc pas décrite par sa trajectoire $x(t)$ comme une particule classique, mais par un champ d’amplitude de probabilité $\psi(x,t)$. Ainsi, une particule quantique ne possède pas de position  bien définie, car la position $x$ est ici une \emph{variable aléatoire} de densité de probabilité $\rho(x,t)=|\psi(x,t)|^2$, voir \Fig{trajectoire}.  À partir de cette densité de probabilité, on peut définir les grandeurs suivantes.
\begin{figure}[H]
  \centering
  \subfigure[]{\label{fig:}
  \includegraphics[width=0.47\textwidth]{x_classique.pdf}
  }
  \subfigure[]{\label{fig:}
  \includegraphics[width=0.47\textwidth]{x_quantique.pdf}
  }
  \caption{
    (a) Trajectoire classique d'une particule (ligne bleue pleine). La trajectoire est la donnée de la position fonction du temps $x(t)$. (b) En mécanique quantique, la trajectoire n'existe pas, mais est remplacée par un champ de densité de probabilité $\rho(x,t)$. Ce qui était la trajectoire classique correspond ici à la position moyenne $\langle x\rangle(t)$ (ligne bleue pleine). Un écart-type $\Delta x(t)$ caractérise la dispersion typique de la variable position $x$.
  }
  \label{fig:trajectoire}
\end{figure}

\begin{liste}
  \item La \imp{position moyenne}  définie par
  \begin{equation}
    \langle x\rangle(t)=\int_{-\infty}^{+\infty}x \rho(x,t)\,\d x,
    \label{eqn:moyenne}
  \end{equation}
La position moyenne $\langle x\rangle(t)$ dépend du temps $t$ en général : c'est cette grandeur qui  s'approche le plus de la notion classique de trajectoire.
% Comment détermine-t-on la distribution de probabilité en pratique ? Pour cela, on commence par préparer $N$ particules dans le même état quantique: elles sont alors décrites par la même fonction d'onde. Puis on fait une mesure, par exemple de la position de ces particules à l'instant $t$. On obtient alors une \important{distribution statistique} des $N$ particules qui permet de déterminer $\rho(x,t)$.

\begin{figure}[h!]
\centering
\includegraphics[width=0.75\textwidth]{distribution.pdf}
\caption{Exemple d'histogramme obtenu lors de la mesure de position pour $N$ particules préparées dans le même état $\psi$. La courbe pointillée rouge représente la distribution de probabilité $\rho(x,t)=\abs{\psi(x,t)}^2$.}
\label{fig:distribution}
\end{figure}
\item L'\imp{indétermination} sur la position, définie par l'écart-type $\Delta x(t)$ de la distribution, donné par 
\begin{equation}
\left(\Delta x\right)^2=\langle x^2\rangle-\langle x\rangle^2 =\int_{-\infty}^{+\infty} x^2\rho(x,t)\mathrm{d}x-\left(\int_{-\infty}^{+\infty} x\rho(x,t)\d x\right)^2.
\label{eqn:variance}
\end{equation}
À un instant $t$ donné, la particule se trouve en moyenne autour de la position $\langle x\rangle(t)$ avec une fluctuation de l'ordre de $\Delta x(t)$. Cette fluctuation est une indétermination intrinsèque de la position, et ne provient d'incertitudes de mesure.
% Cette grandeur $\Delta x$ représente l'\imp{indétermination} sur la position $x$ de la particule.
%D'un point de vue statistique (si on réalise un grand nombre de fois l'expérience pour le même état initial) on trouvera les particules positionnées autour de $\langle x\rangle(t)$ avec une dispersion typique des points de mesure de l'ordre de $\Delta x(t)$.

\end{liste}

Les fluctuations statistiques autour de la moyenne de certaines grandeurs, comme la position $x$ et la quantité de mouvement $p_x$, ne sont pas indépendantes. 
% Ce résultat important est admis dans le cadre du programme.
\begin{theorem}[Relation d'indétermination de Heisenberg]\label{form:}
  Les indéterminations sur la position $\Delta x$ et sur la quantité de mouvement $\Delta p_x$ vérifient
  $$
  \Delta x \Delta p_x \geq \frac{\hbar}{2}.
  \label{eqn:indetermination}
  $$
\end{theorem}
% où $\hbar=\frac{h}{2\pi}\simeq\SI{1,05e-34}{J.s}$ est la constante de Planck réduite et $\Delta p_x$ désigne l'indétermination sur la quantité de mouvement (selon $\vv{e_x}$).
Le même genre d'inégalité s'obtient pour les autres directions d'espace ($y$ et $z$). Cette inégalité signifie que l'indétermination sur la position $x$ et l'impulsion $p_x$ ne peuvent pas être arbitrairement petite : mieux l'impulsion est définie ($\Delta p_x$ petit), et moins bien la position est définie ($\Delta x$ grand). La notion de trajectoire perd son sens en mécanique quantique. L'inégalité d'Heisenberg n'imposerait aucune condition dans la limite hypothétique où $\hbar\to 0$ (limite classique).

\begin{remarque}
Il est préférable d'utiliser le terme indétermination plutôt qu'incertitude, car l'indétermination est ici fondamentale (intrinsèque à la physique) et n'a rien à voir l'imprécision due à la mesure. Même dans le cas limite où nous aurions des appareils de précision infinie, l'incertitude classique serait nulle mais l'indétermination quantique subsisterait d'après l'inégalité de Heisenberg.
\end{remarque}

\section{Équation de Schrödinger}
\subsection{Énoncé}
\label{subsec:schrodinger}
De même que la vibration lumineuse est décrite par une équation d’onde réelle (l’équation de d’Alembert), la fonction d’onde $\psi(x,t)$ est décrite par une équation d'onde complexe appelée \imp{équation de Schrödinger}. Il s'agit d'un postulat de la mécanique quantique ondulatoire (au même titre que les équations de Maxwell en électromagnétisme ou la deuxième loi de Newton en mécanique classique).
\begin{theorem}[Équation de Schrödinger]\label{form:}
La fonction d'onde $\psi(x,t)$ d'une particule de masse $m$ d'énergie potentielle $V(x,t)$  vérifie
\begin{equation*}
  i\hbar\frac{\partial\psi}{\partial t}=-\frac{\hbar^2}{2m}\frac{\partial^2\psi}{\partial x^2}+V\psi.
\end{equation*}
\end{theorem}
  % Cette équation de Schrödinger décrit l'évolution temporelle de toute particule quantique massive non relativiste.
% \begin{remarque}
  L'évolution de la fonction d'onde est donc déterministe\footnote{Ce qui n'empêche évidemment pas la théorie d'être probabiliste, comme nous l'avons vu en Sec.~\ref{sec:description_proba}, car ce n'est pas parce que la fonction d'onde est connue que les grandeurs qui décrivent la particule (par exemple la position) le sont.}.
L’énergie potentielle $V(x,t)$ est appelée plus simplement \emph{potentiel} dans ce contexte.
On rappelle que la constante $\hbar=\frac{h}{2\pi}\simeq\SI{1,05e-34}{J.s}$ désigne la constante de Planck réduite.
% En dimensions supérieures, le terme $\frac{\partial^2\psi}{\partial x^2}$ de l'équation devient $\Delta\psi$ (Laplacien scalaire).
% \end{remarque}
Quelques remarques physiques sur l'équation de Schrödinger.
\begin{liste}
  \item Elle est bien homogène, puisque $\hbar \dr{}{t}$, $\frac{\hbar^2}{2m} \ddr{}{x}$ et $V$ ont la dimension d’une énergie.
\item Il s'agit d'une équation \imp{linéaire}, ce qui assure le principe de superposition des solutions.
%, et permet d'expliquer les phénomènes d'interférences des ondes de matière mentionnés en Sec.~\ref{subsec:interfs}.
\item Elle fait intervenir une dérivée première en $t$ : il suffit de connaître la fonction d'onde initiale $\psi(x,t=0)$ pour connaître  $\psi(x,t)$ à tout instant. En cela, l'équation de Schrödinger diffère de la plupart des équations d'onde, qui sont des équations aux dérivées partielles du second ordre en temps (d'Alembert et équations proches). Elle se rapproche d'une équation de diffusion. 
\end{liste}


\subsection{États stationnaires}

\parag{Forme d’un état stationnaire}
On suppose pour simplifier que le potentiel ne dépend pas du temps, $V(x,t)=V(x)$. On peut alors chercher des solutions de l'équation de Schrödinger sous la forme d'ondes stationnaires, c’est-à-dire d’ondes séparables en les variables $x$ et $t$,
% \begin{equation}
  \eqbox[]{
    % \label{eq:OS_schro}
    \psi(x,t)=\varphi(x)g(t),
    }
% \end{equation}
où $\varphi(x)$ et $g(t)$ sont deux fonctions complexes.
% \begin{definition}[État stationnaire]\label{def:}
%   On appelle \imp{état stationnaire} d'une particule évoluant dans un potentiel $V(x)$ indépendant du temps une fonction d'onde séparable en les variables $x$ et $t$,
%   $$
%   \psi(x,t)=\varphi(x)g(t).
%   $$
% \end{definition}
%
En imposant la normalisation de la fonction d'onde,
\begin{equation}
|g(t)|^2\underbrace{\int_{-\infty}^{\infty}|\varphi(x)|^2\,\d x}_{\rm constante}=1.
\end{equation}
on voit que $|g(t)|$ est une constante, qu'on peut prendre égale à $1$ sans perte de généralité. Il suffit de normaliser $\varphi(x)$ séparément. Ainsi $g(t)=\e^{i\alpha(t)}$ avec $\alpha(t)$ une fonction réelle du temps. Un état stationnaire d'un système quantique s'écrit sous la forme
% \begin{theorem}[État stationnaire]\label{form:stat}
  La fonction d’onde d’un état stationnaire est de la forme
  % \begin{equation*}
    \eqbox[]{
      \psi(x,t)=\varphi(x)\,\e^{i\alpha(t)},\quad \text{avec} \quad
      \int_{-\infty}^{+\infty}|\varphi(x)|^2\,\d x = 1.
      }
  % \end{equation*}
% \end{theorem}


%   On en tire donc que pour que la fonction d'onde soit normalisée, il faut que
% \begin{equation}
% \int_{x\in\mathcal{D}}|\varphi(x,t)|^2\,\d x=1.
% \label{eqn:norm2}
% \end{equation}


Dans un état stationnaire, la densité de probabilité de présence de la particule $\rho(x)=|\varphi(x)|^2$ est indépendante du temps. Par conséquent, toutes les grandeurs moyennes et indéterminations sont indépendantes du temps ($\moy{x},\Delta x$, etc.).
%\item le courant de probabilité est aussi indépendant du temps, et s'écrit
%\begin{equation}
%\vec{J}(\vec{r})=\frac{-i\hbar}{2m}\left[\varphi^*(\vec{r})\grad{\varphi}(\vec{r})-\varphi(\vec{r})\grad{\varphi^*}(\vec{r})\right].
%\end{equation}
% \end{itemize}

\parag{Solution stationnaire de l'équation de Schrödinger}

On cherche maintenant une solution de la forme de la \Formu{stat} à l'équation de Schrödinger. En substituant il vient
\begin{equation}
-\hbar\,\varphi\frac{\d\alpha}{\d t}\e^{i\alpha}=-\frac{\hbar^2}{2m}\ddd{\vphi}{x}\e^{i\alpha}+V\varphi\e^{i\alpha}
\end{equation}
qui se simplifie en
\begin{equation}
-\hbar\frac{\d\alpha}{\d t} =-\frac{\hbar^2}{2m}\frac{1}{\varphi}\ddd{\vphi}{x}+V.
\end{equation}
Le membre de gauche ne dépend que du temps $t$, tandis que celui de droite ne dépend que de l'espace $x$, donc ces deux quantités sont égales à une constante réelle $E\geq 0$ qui par analyse dimensionnelle a la dimension d'une énergie. L'équation temporelle se réécrit donc
%\footnote{À une constante additive près, qu'on peut prendre nulle en redéfinissant l'origine des temps.}
\begin{equation}
-\hbar\dd{\alpha}{t}=E \qquad\Rightarrow\quad \alpha(t)=-\frac{Et}{\hbar},
\end{equation}
où l’on a choisit l’origine des temps de sorte que $\alpha(t=0)=0$.

% \begin{prop}[État stationnaire de l'équation de Schrödinger]\label{prop:}
%   Les états stationnaires de l'équation de Schrödinger sont les états d'énergie $E$ fixée, de la forme
%   \begin{equation*}
%   \psi(x,t)=\varphi(x)\,\e^{-i\frac{Et}{\hbar}}, \qqavec
%   \int_{-\infty}^{+\infty}|\varphi(x)|^2\,\d x = 1.
%   \label{eqn:etat_stat3}
% \end{equation*}
% \end{prop}
La fonction d’onde d’un état stationnaire oscille dans le temps à une pulsation $\omega$ que vérifie la relation  de Planck-Einstein donnée dans le \Tab{dico_Ep},
$
  E=\hbar\omega.
$
% Cela justifie la $E=\hbar \omega$ }.
% où désigne l'énergie totale de la particule\footnote{Attention donc, cette relation $E=\hbar\omega$ n'est valable que pour des états \important{stationnaires}.},
De plus, $\varphi(x)$ vérifie l’équation  suivante.

\begin{prop}[État stationnaire]\label{form:schro_stat}
  Les états stationnaires de l'équation de Schrödinger sont les états d'énergie $E$ fixée, de fonction d'onde
  \begin{equation*}
  \psi(x,t)=\varphi(x)\,\e^{-i\frac{Et}{\hbar}}.
  % \qqavec
  % \int_{-\infty}^{+\infty}|\varphi(x)|^2\,\d x = 1.
  \label{eqn:etat_stat3}
\end{equation*}
\end{prop}

\begin{theorem}[Équation de Schrödinger indépendante du temps]\label{form:schro_stat}
La fonction d'onde spatiale $\vphi(x)$ vérifie l'équation de Schrödinger indépendante du temps
  \begin{equation*}
  -\frac{\hbar^2}{2m}\frac{\d^2\varphi}{\d x^2}(x)+V(x)\varphi(x)=E\varphi(x).
\end{equation*}
\end{theorem}

% Les résultats \eqref{eqn:etat_stat3} et \eqref{eqn:schro_stat} sont absolument à retenir.

\parag{Conditions aux limites}
L'équation de Schrödinger indépendante du temps \Formu{schro_stat} admet un nombre infini de solutions pour la fonction d'onde spatiale $\varphi(x)$. Le choix de la bonne solution pour un problème donné dépend des conditions aux limites. On admet les propriétés suivantes.
\begin{enumerate}[label=\boxed{{\rm CL}\arabic*},leftmargin=40pt]
\item Toute fonction d'onde normalisable doit vérifier
\eqbox[]{
\lim\limits_{x\to\pm\infty}\varphi(x)=0.
}
\begin{remarque}
  Cette contrainte ne s'applique pas aux OPPH, qui sont des modèles d'onde non réalistes.
\end{remarque}
\item La fonction d'onde est $\mc{C}^1$ par morceaux. Plus spécifiquement,
\eqbox[]{
\varphi(x)\text{ est continue}
} 
en tout point $x\in \R$, et
\eqbox[]{
  \vphi'(x)\text{ est continue}
}
en tout point $x\in \R$ où $V(x)$ est continu ou subit une discontinuité finie. Cette propriété de continuité sera donc valable pour le puits de potentiel fini mais pas le puits de potentiel infini.
\end{enumerate}
% \vspace{-2mm}
% \begin{remarque}
% Nous verrons que les solutions physiquement acceptables qui satisfont ces conditions correspondent à des valeurs discrètes de l'énergie $E$. On dira qu'il y a quantification de l'énergie et donc des modes propres, par analogie avec la quantification par les conditions aux limites déjà observée sur la corde de Melde en physique des ondes.
% \end{remarque}
\subsection{Évolution d'une particule libre}
On s'intéresse maintenant au états stationnaires d’une particule \imp{libre}, c'est-à-dire évoluant dans un potentiel $V(x,t)$ constant, qu’on peut prendre nul : $V=0$. L’équation de Schrödinger indépendante du temps se simplifie en
\begin{equation*}
-\frac{\hbar^2}{2m}\frac{\d^2\varphi}{\d x^2}=E\varphi.
\end{equation*}
% On suppose que la particule ne peut explorer qu’un segment de longueur $L$ de l’espace. La fonction d’onde n’est donc non nul que pour $0\leq x \leq L$.
\begin{definition}[Particule libre]\label{def:}
  On appelle \imp{particule libre} une particule évoluant dans un potentiel nul, $V(x,t)=0$.
\end{definition}

%  qu'on peut prendre nul par translation de l'origine des énergies potentielles. La fonction d'onde $\psi(x,t)$ vérifie alors l'équation de Schrödinger suivante
% \begin{equation}
% i\hbar\frac{\partial \psi}{\partial t}=-\frac{\hbar^2}{2m}\frac{\partial^2\psi}{\partial x^2}.
% \label{eqn:schro_libre}
% \end{equation}
%
% \subsubsection{De l'onde de de Broglie à l'équation de Schrödinger}
% Dans ce bref paragraphe, nous montrons qu'il est possible de retrouver de façon heuristique l'équation de Schrödinger simplement à partir des arguments de de Broglie. On cherche à décrire le comportement des particules sous forme ondulatoire, donc on cherche $\psi(x,t)=\frac{1}{\sqrt{L}}\e^{-i(\omega t-kx)}$ où la constante est fixée par normalisation sur le domaine $\mc{D}$ de longueur $L$. La relation de de Broglie $\lambda=h/p$ s'écrit de façon équivalente $k=p/\hbar$. Pour relier $E$ et $\omega$, de Broglie postule comme pour le photon que la relation se met sous la forme $E=\hbar\omega$ de sorte que
% \begin{equation}
% \psi(x,t)=\frac{1}{\sqrt{L}}\e^{-\frac{i}{\hbar}(Et-px)}.
% \end{equation}
% Mais pour des particules libres (absence d'énergie potentielle extérieure) on a de plus $E=E_{\rm{cin}}=\frac{1}{2}mv^2=\frac{p^2}{2m}$ où $p=mv$ désigne la quantité de mouvement. On peut alors former une équation d'onde, car
% \begin{equation}
% \begin{cases}
% \frac{\partial\psi}{\partial t}=-\frac{i}{\hbar}E\psi=-\frac{i}{\hbar}\frac{p^2}{2m},\\
% \frac{\partial^2\psi}{\partial x^2}=\frac{-p^2}{\hbar^2}\psi.
% \end{cases}
% \end{equation}
% Autrement dit l'équation d'onde d'une particule libre est de la forme
% \begin{equation}
% i\hbar\frac{\partial\psi}{\partial t}=-\frac{\hbar^2}{2m}\frac{\partial^2\psi}{\partial x^2},
% \end{equation}
% qui est effectivement l'équation de Schrödinger en l'absence de potentiel extérieur.

\parag{Fonction d'onde} On peut récrire l’équation de Schrödinger sous la forme
% Nous suivons la démarche inverse de celle développée précédemment: on postule l'équation de Schrödinger et nous allons en déduire la relation de dispersion des ondes de matière.
\begin{equation*}
\frac{\d^2\varphi}{\d x^2} + \frac{2mE}{\hbar^2} \varphi=0.
\end{equation*}
Pour obtenir des solutions oscillantes et pas des solutions qui décroissent exponentiellement, on doit imposer $E>0$. On peut alors poser $k^2 = \frac{2mE}{\hbar^2}$ avec $k>0$ le module du vecteur d’onde. La fonction d’onde est donc de la forme
\begin{equation}
\psi(x,t)=A\e^{-i(\omega t-kx)} +B\e^{-i(\omega t+kx)}.
\label{eqn:oppqh}
\end{equation}
Il s’agit d’une solution sous forme d’OPPH ; le premier terme décrit une OPPH qui se propage suivant les $x$ croissants, et le second terme une OPPH qui se propage suivant les $x$ décroissants. Une OPPH suivant les $x$ croissants s'écrit donc 
\begin{equation}
\psi(x,t)=A\e^{-i(\omega t-kx)}.
\end{equation}
\begin{liste}
  \item Si la particule évolue dans une région infini de l'espace, la constante $A$ est déterminée. C'est le cas d'une OPPH incidente sur un puits ou une barrière de potentiel.
  \item Si la particule évolue dans une boîte de longueur $L$, la constante $A$ est déterminée par la condition de normalisation,
  \begin{equation}
    \int_{0}^{L} |\psi(x,t)|^2 \d x = |A|^2 L = 1 \Rightarrow |A| = \frac{1}{\sqrt{L}}.
  \end{equation}
  Sans perte de généralité, on peut supposer $A$ réel positif, et donc $A=\frac{1}{\sqrt{L}}$.
\end{liste}
\begin{prop}[États stationnaires d’une particule libre]\label{prop:}
  Les états stationnaires d’une particule libre sont des OPPH complexe, de la forme
% On cherche les solutions sous la forme d'une OPPH
\begin{equation*}
\psi(x,t)=\frac{1}{\sqrt{L}}\,\e^{-i(\omega t-kx)},
% \label{eqn:oppqh}
\end{equation*}
pour une onde se propageant suivant les $x$ croissants dans une boîte de longueur $L$.
\end{prop}
% En substituant dans l'équation de Schrödinger libre, on obtient qu'une OPPH est solution si

\begin{remarque}
  IMPORTANT : un état stationnaire en mécanique quantique est différent d'une onde stationnaire en physique classique. En effet, un état de la forme $\psi(x,t)=\psi_0\,\e^{-i(\omega t-kx)}$ est un état stationnaire, parce que les exponentielles spatiale et temporelle peuvent se factoriser, mais la représentation complexe d'une vibration lumineuse, par exemple, $\udl{s}(x,t)=s_0\,\e^{-i(\omega t-kx)}$, n'est pas une onde stationnaire, puisqu'elle représente l'onde réelle $s(x,t)=s_0\,\cos(\omega t-kx)$, qui n'est pas factorisable. En conclusion, une OPPH représente un état stationnaire en mécanique quantique mais ne représente pas une onde stationnaire en physique classique.
\end{remarque}

\parag{Relation de dispersion}
La relation de dispersion $\omega(k)$ s’obtient à partir de la définition de $k$ et de la relation de Planck-Einstein :
\begin{equation}
  E = \frac{\hbar^2 k^2}{2m} = \hbar\omega.
\end{equation}
\begin{remarque}
  On voit au passage que si $p=\hbar k$ est l’impulsion de la particule, alors $E = \frac{p^2}{2m}$ correspond à l’énergie cinétique de la particule. Pour une particule classique, on aurait en effet $p=mv$ et donc $E = \frac{1}{2} m v^2$.
\end{remarque}
\begin{theorem}[Relation de dispersion]\label{form:}
  La relation de dispersion de l'équation d'onde d'une particule libre est
  % Pour être solution de l’équation de Schrödinger, la fonction d’onde stationnaire d’une particule libre doit vérifier la relation de dispersion
  \begin{equation*}
  \omega=\frac{\hbar k^2}{2m}.
\end{equation*}
\end{theorem}
Contrairement au photon (particule quantique associée aux ondes électromagnétiques) pour lequel $\omega=ck$, les particules massives ont une relation de dispersion \imp{non linéaire}. Au cours de la propagation, il y a donc \important{dispersion} du paquet d'onde. En revanche, $k\in\mathbb{R}$ et donc le paquet d'onde se propage sans atténuation.
%, comme on pouvait s'y attendre en l'absence de force extérieure. En utilisant la relation de Planck-Einstein $E=\hbar\omega$ et celle de de Broglie $p=\hbar k$ on retrouve alors ce que nous avons dit précédemment à savoir $E=\frac{p^2}{2m}$ (énergie cinétique).

\parag{Solution en paquet d'onde}
Une OPPH dans l’espace entier ($L\rightarrow \infty$) serait d'extension spatiale infinie, et n’est pas normalisable. Une telle solution n’est pas physique, et doit être remplacé par un paquet d'onde
\begin{equation}
\label{eq:paquet}
 \psi(x,t)=\int_{\mathbb{R}} \d k \,A(k)\,\e^{-i[\omega(k)t-kx]}.
 \end{equation}
%  \begin{remarque}
% La relation précédente est une transformée de Fourier: $a(k)$ représente la contribution de l'onde caractérisée par le nombre d'onde $k$ et la fréquence $\omega(k)=\frac{\hbar k^2}{2m}$ (d'après la relation de dispersion) à la fonction d'onde totale. Autrement dit, on décompose un paquet d'onde en une somme d'harmoniques.
% \end{remarque}

Supposons le paquet d'onde centré autour d'un vecteur d’onde de module $k_0$ et d’une pulsation $\omega_0 = \omega(k_0)$. L’amplitude spectrale $A(k)$ ne prend des valeurs significatives que dans un intervalle de largeur $\Delta k$ centré autour de $k_0$. On peut alors développer l'argument de l'exponentielle dans l'Éq.~\eqref{eq:paquet} au premier ordre en $k-k_0$ (approximation de faible dispersion) pour trouver
\begin{equation}
\begin{aligned}
\omega(k)t - k x &=\left[\omega_0+\left(k-k_0\right)\dd{\omega}{k}(k_0)\right]t - k_0 x - \left(k-k_0\right)x ,\\
&= k_0(v_\vphi t- x) + \left(k-k_0\right)\left(v_\mathrm{g} t-x\right),
\end{aligned}
\end{equation}
où l’on a introduit les vitesses $v_\vphi$ et $v_g$.
\begin{liste}
   \item La \imp{vitesse de groupe} a pour expression
  \begin{equation}
    v_g = \dd{\omega}{k}(k_0) = \frac{\hbar k_0}{m} = \frac{p}{m}
  \end{equation}
  où $p=\hbar k_0$ désigne l’impulsion moyenne de la particule. On voit que $p=mv_g$, ce qui signifie que la vitesse de groupe est l’équivalent de la vitesse $v$ de la particule en physique classique.
  \item La \imp{vitesse de phase} a pour expression
  \begin{equation}
    v_\vphi = \frac{\omega_0}{k_0} = \frac{\hbar k_0}{2m} = \frac{v_g}{2}.
  \end{equation}
  La vitesse de phase est deux fois plus faible que la vitesse de groupe.
\end{liste}
Le paquet d'onde s'écrit donc sous la forme, après changement de variable $k \mapsto k+k_0$,
\begin{equation}
\psi(x,t)=\underbrace{\e^{-i\left(\omega_0 t - k_0 x\right)}}_{\text{OPPH qui se propage à~} v_{\varphi}}\times\underbrace{\int_{\mathbb{R}}\d k \, A(k_0+ k)\,\e^{-i k\left(v_\mathrm{g}t - x\right)}}_{\text{enveloppe qui se propage à~} v_{\rm{g}}},
\end{equation}

\begin{figure}[H]
  \centering
  \includegraphics[width=0.9\textwidth]{paquet_onde.jpg}
  \caption{(a) Module de la fonction d'onde, et donc de l'enveloppe d'un paquet d'onde. (b) Module de l'amplitude spectrale du paquet d'onde.}
  \label{fig:paquet}
\end{figure}

% La fonction d'onde de la particule peut alors se réécrire comme le produit d'un terme exponentiel oscillant qui se propage à la vitesse de phase $v_\varphi(k_0)$ et d'une enveloppe $\mathcal{A}(x-v_\mathrm{g}(k_0)t)$ qui se propage à la vitesse de groupe $v_\mathrm{g}(k_0)$ sans se déformer, où
% \begin{equation}
% \mathcal{A}(u)=\int_{\delta k\in\mathbb{R}}\frac{\mathrm{d}\delta k}{2\pi}\ a(k_0+\delta k)\,\e^{i\delta k u}.
% \end{equation}
L’enveloppe est d’amplitude maximale lorsque $x=v_\mathrm{g} t$. Le pic du paquet d'onde décrit donc la trajectoire classique $x(t)=vt$ de la particule. De plus, d'après les propriétés de la transformée de Fourier (TF), les étendues spatiale $\Delta x$ et spectrale $\Delta k$ de l'enveloppe vérifient
\begin{equation}
  \Delta x \Delta k \geq \frac{1}{2}.
\end{equation}
Or l'impulsion vérifier la relation de de Broglie $p=\hbar k$, et donc
\begin{equation}
  \Delta x \Delta p \geq \frac{\hbar}{2}.
\end{equation}
On retrouve la relation d'indétermination de Heisenberg, qui est donc une conséquence directe de la description de particules en terme de fonction d'onde.

\subsection{Courant de probabilité}
Une conséquence de l'équation de Schrödinger est l'équation de conservation de la probabilité. On peut montrer (démonstration hors programme) que cette équation prend la forme suivante.

\begin{theorem}[Équation de conservation de la probabilité]\label{form:}
  À une dimension, l'équation de conservation de la probabilité s'écrit
  \begin{equation*}
    \frac{\partial\rho}{\partial t}+\frac{\partial j}{\partial x}=0,
  \end{equation*}
  où $\rho=\abs{\psi}^2$ est la densité de probabilité et $j=\rho v_{\rm{g}} = \rho \frac{\hbar k}{m}$ est la densité de courant de probabilité.
\end{theorem}
\begin{remarque}
  
  On notera que ce courant s'interprète comme d'habitude comme le produit de la densité $\rho$ par la vitesse $v_g$.
  À 3D, l'équation de conservation devient $\frac{\partial\rho}{\partial t}+\dive{\j}=0$ avec $\j = \rho \,\Vec{v_g}$.
\end{remarque} 
  
\begin{exemple}{Démonstration de l'équation de conservation de la probabilité}
  On commence par réécrire l'équation de Schrödinger libre et son complexe conjugué
\begin{equation*}
\begin{cases}
i\hbar\frac{\partial\psi}{\partial t}=-\frac{\hbar^2}{2m}\frac{\partial^2\psi}{\partial x^2}\\[1mm]
-i\hbar\frac{\partial\psi^{\star}}{\partial t}=-\frac{\hbar^2}{2m}\frac{\partial^2\psi^{\star}}{\partial x^2}\\
\end{cases}
\end{equation*}
\begin{quest}
  \item	Sachant que $\rho=\psi\psi^{\star}$, déterminer $\dr{\rho}{t}$ à partir des équations précédentes.
  \item Définir et simplifier l'expression de la densité de courant de probabilité $j$, et démontrer l'équation de conservation de la probabilité.
\end{quest}

\examplesol{
  \begin{quest}
    \item	Or $\rho=\psi\psi^{\star}$ et donc
    \begin{equation*}
    \frac{\partial\rho}{\partial t}=\psi\frac{\partial\psi^\star}{\partial t}+\psi^\star\frac{\partial\psi}{\partial t}=\frac{-i\hbar}{2m}\left(\psi\frac{\partial^2\psi^{\star}}{\partial x^2}-\psi^{\star}\frac{\partial^2\psi}{\partial x^2}\right)
    \end{equation*}
    \eqbox{
      \frac{\partial\rho}{\partial t}=\frac{-i\hbar}{2m}\frac{\partial}{\partial x}\left(\psi\frac{\partial\psi^{\star}}{\partial x}-\psi^{\star}\frac{\partial\psi}{\partial x}\right).
    }
    \item 
    Autrement dit le courant de probabilité $j_x$ s'écrit
    \begin{equation*}
    j=\frac{i\hbar}{2m}\left(\psi\frac{\partial\psi^{\star}}{\partial x}-\psi^{\star}\frac{\partial\psi}{\partial x}\right).
    \end{equation*}
    En utilisant $\psi(x,t)=\frac{1}{\sqrt{L}}\,\e^{-i(\omega t-kx)}$ et donc $\psi^{\star}(x,t)=\frac{1}{\sqrt{L}}\,\e^{i(\omega t-kx)}$ il vient
    \begin{equation*}
    j=\frac{i\hbar}{2m}(-2ik)\abs{\psi}^2=\frac{\hbar k}{m}\abs{\psi}^2,
    \end{equation*}
    soit 
    \eqbox{
      j=\rho v_{\rm{g}}.
    }
    La densité de courant de probabilité à 3 dimensions s'écrit de façon générale 
    \begin{equation*}
      \j = \frac{\hbar}{2mi}\left(\psi^{\star}\grad\psi-\psi\,\grad\psi^{\star}\right) = \frac{\hbar}{m} \im \left(\psi^{\star}\grad\psi\right).
    \end{equation*}
  \end{quest}
}
\end{exemple}

\subsection{Lien avec les orbitales atomiques}
 Il est erroné de visualiser un atome comme un ensemble d’électrons en orbite autour d’un noyau. Cette vision classique doit être remplacée par une description quantique pour expliquer la plupart des propriétés atomiques. Un électron de charge $-e$ de masse $m$ subit la force d'attraction coulombienne du noyau atomique de charge $+Ze$, de potentiel 
 \begin{equation}
    V(r) = -\frac{Ze^2}{4\pi\epz r}.
 \end{equation}
Ce problème de mécanique 3D en coordonnées sphériques $(r,\theta,\vphi)$ se réduit à un problème de mécanique 1D sur la variable $r \in [0,+\infty[$, de potentiel effectif 
  % \begin{equation}
  \eqbox[Veff]{
    V_{\rm{eff}}(r) = -\frac{Ze^2}{4\pi\epz r} + \frac{L^2}{2mr^2},
    }
  % \end{equation}
  où $\Vec{L}$ désigne le moment cinétique orbital de l'électron. Nous avons vu les règles de quantification du moment cinétique orbital en mécanique quantique,
  \begin{equation}
    \begin{cases}
      L^2 & = \hbar^2 l(l+1),\\ 
      L_z & = \hbar m_l,
    \end{cases}
  \end{equation}
  % \begin{equation}
    où $\ell$ (entier) est le nombre quantique azimutal et $m_\ell$ (entier) le nombre quantique magnétique. La résolution de l'équation de Schrödinger montre que les électrons sont décrits par des fonctions d’ondes stationnaires appelées \imp{orbitales atomiques}, de la forme
  \eqbox[]{
    \psi_{n, \ell, m_\ell}(M) = R_{n,\ell}(r) Y_{\rm orb}(\theta,\vphi),
    }
% \end{equation}
où $n=1,2,3,...$ est le nombre quantique principal et \guill{orb} est déterminée par les nombres quantiques $\ell=0,1,2,...,n-1$ et $m_\ell=-\ell,-\ell-1,...,\ell-1,\ell$, et représente les sous-couches $s$, $p_x,p_y,p_z$, etc, voir \Tab{orbitals}. À $n$ donné, le nombre quantique azimutal $\ell$ ne peut donc prendre que $n$ valeurs ; à $\ell$ donné, le nombre quantique magnétique $m_\ell$ peut prendre $2\ell+1$ valeurs. 

\begin{table}[h!]
\centering
\begin{tabular}{|c|c|c|c|}
\hline
$\ell$ & 0 & 1 & 2 \\
\hline
Couche atomique & 
$s$ & $p$ & $d$  \\ 
\hline
Nombre de sous-couches ($2\ell+1$) & 1 & 3 & 5 \\
\hline
Sous-couches & $s$ & $p_x$, $p_y$, $p_z$ & $d_{xy}$, $d_{yz}$, $d_{xz}$, $d_{x^2-y^2}$, $d_{z^2}$ \\
\hline
\end{tabular}
\caption{
  Les sous-couches atomiques pour les premiers nombres quantiques azimutaux $\ell$. Les sous-couches sont représentées sur la \Fig{orbitales}.
}
\label{tab:orbitals}
\end{table}


% Par exemple, l’orbitale $2p_x$ est décrite par la fonction d’onde $\psi_{2p_x}(M) = R_2(x) Y_{p_x}(\theta,\vphi)$. La fonction d’onde est séparable dans les variables de distance $r$ et d'angle $\theta,\vphi$.
\begin{liste}
  \item La fonction radiale $R_n(r)$ est réelle et ne dépend que de la distance au noyau $r$. Elle décroît exponentiellement avec la distance sur une distance typique $na_0$ où $a_0 \simeq \SI{0.5}{\angstrom}$ est le rayon de Bohr. Elle possède $n-\ell-1$ annulations, ou nœuds, comme illustré sur la \Fig{orbitale_s} Le nombre quantique principal $n$ détermine donc sa distance moyenne au noyau, ainsi que son énergie  
  \begin{equation}
    E_n = -\frac{Z^2 E_1}{n^2}, \text{ où } E_1 = \frac{e^4 m}{32\pi^2\epz^2\hbar^2} \simeq \SI{13.6}{\eV}
  \end{equation}
  est l'énergie de Rydberg.
  % Les premières fonctions pour $n=1,2,3$ sont représentées sur la \Fig{}.
   Ce n’est pas cette partie de la fonction d’onde qui intéresse les chimistes, parce qu’elle ne détermine pas dans les liaisons covalentes, et donc intervient peu dans les propriétés chimiques des éléments.
   \begin{figure}[H]
    \centering
    \includegraphics{orbitale_s.pdf}
    \caption{Premières orbitales atomiques de type $ns$, avec ici $n=1,2,3$. Comme $\ell=0$ pour une orbitale $s$, les orbitales $1s$, $2s$ et $3s$ possèdent respectivement $0$, $1$ et $2$ n\oe{}uds.}
    \label{fig:orbitale_s}
   \end{figure}
  \item La fonction $ Y_{\rm orb}(\theta,\vphi)$ dépend de la direction dans l’espace, repérée par les angles $\theta$ et $\vphi$. C’est cette fonction qui dicte l’orientation des liaisons covalentes dans une molécule. Ces fonctions s’appellent \imp{harmoniques sphériques}. En pratique on choisit une base d’harmoniques sphériques \emph{réelles}, de sorte que $\psi(M)$ soit aussi réel.
\end{liste}
Quelques orbitales atomiques sont représentées sur la \Fig{orbitales}. Pour représenter les orbitales, on représente la surface qui contient l’ensemble des points $M(r,\theta,\vphi)$ telle que $r=\cte \times| Y_{\rm orb}(\theta,\vphi)|$. Il s’agit donc d’une surface \guill{iso-probable} : tous les points $M$ de cette surface correspondent à une même densité de probabilité $\rho(M)$. On choisit la constante de sorte que l’électron ait une certaine probabilité de présence d’être à l’intérieur de la surface, par exemple $\SI{90}{\%}$. La couleur de la surface (ici orange ou bleu) indique le signe de la fonction d’onde (positive ou négative).
% Nous l'avons mentionné, une description probabiliste est nécessaire. Par analogie avec les phénomènes ondulatoires, nous allons décrire l'état d'une particule par une amplitude de vibration scalaire.

\begin{figure}[h!]
\centering
\includegraphics[width=0.7\textwidth]{orbitales.pdf}
\caption{Représentation des premières orbitales réelles pour un atome hydrogénoïde. La couleur représente le signe de la fonction d'onde $\psi_{n,\rm orb}(r,\theta,\vphi)$.
% Les n\oe{}uds correspondent à une annulation de la probabilité de présence de l'électron.
}
\label{fig:orbitales}
\end{figure}


\begin{exemple}{Propriétés des orbitales $p$}
  Les fonctions d’onde en coordonnées sphériques des orbitales de type $2p$ de l'atome d'hydrogène sont de la forme $\psi_{2p}(r,\theta,\vphi) = R_{2,2}(r) Y_{p}(\theta,\vphi)$, avec $R_{2,2}(x) =  \frac{1}{\sqrt{24a_0^3}} \frac{r}{a_0}\e^{-r/2a_0}$ et 
  \begin{equation*}
    \begin{cases} 
      Y_{p_x}(\theta,\vphi) = \sqrt{\frac{3}{4\pi}} \sin \theta  \cos \vphi, \\
      Y_{p_y}(\theta,\vphi) = \sqrt{\frac{3}{4\pi}} \sin \theta  \sin \vphi, \\
      Y_{p_z}(\theta,\vphi) = \sqrt{\frac{3}{4\pi}} \cos \theta.
    \end{cases} 
  \end{equation*}

  \begin{quest}
    \item	Quelles sont les valeurs autorisées de $\ell$ au niveau $n=2$ ? À quelle valeur de $\ell$ correspond la couche $2p$ ? Justifier qu'il y a $3$ sous-couches, $2p_x$, $2p_y$ et $2p_z$.
    \item Évaluer l'énergie d'un électron d'un atome d'hydrogène sur un niveau $n=2$. 
    % Montrer que la valeur moyenne de la distance de l'électron au noyau vaut $\moy{r} =  2a_0$ sachant que les fonctions $R_{n,\ell}(r)$ et $Y_{\rm orb}(\theta,\vphi)$ sont normalisées séparément.
    \item Pour visualiser la dépendance angulaire des harmoniques sphériques $Y_{\rm orb}(\theta,\vphi)$, on représente la surface d'iso-probabilité définie par la surface d'équation en coordonnées sphériques  $r = \cte \times |Y_{\rm orb}(\theta,\vphi)|^2$ en coordonnées sphériques. Déterminer l'équation de cette surface pour les orbitales $2p_x$, $2p_y$ et $2p_z$. Pour quelle valeur de $\theta$ et $\vphi$ la probabilité de présence s'annule-t-elle pour chaque orbitale ? En déduire l'existence de plans nodaux.
  \end{quest}
  % $Y_{p}(\theta,\vphi) = \sqrt{\frac{3}{4\pi}} \sin \theta  \sin \vphi$. On représente donc l’orbitale $2p$ par la surface d’équation $r = \cte \times |\sin \theta  \sin \vphi|$ avec $0\leq \theta \leq \pi$ et $0\leq \vphi \leq 2\pi$. On voit que le rayon de cette surface s’annule dans le plan d’équation $\vphi = 0$ et $\vphi = \pi$, soit $y=0$ et $x=0$. Un tel plan est appelé plan nodal, et est un plan d’antisymétrie de la fonction d’onde. Dans le plan d’équation $\theta = \frac{\pi}{2}$, soit $z=0$, on peut représenter plus simplement l’orbitale atomique dans un plan par la courbe d’équation $r = \cte  \times |\sin\vphi|$ en coordonnées polaires. On observe deux lobes dans la direction $(Oy)$.
  
  % avec $Y_{p_x}(\theta,\vphi) = \sqrt{\frac{3}{4\pi}} \sin \theta  \cos \vphi$. On représente donc l’orbitale $2p_x$ par la surface d’équation $r = \cte \times |\sin \theta  \cos \vphi|$ avec $0\leq \theta \leq$
  % et $0\leq \vphi \leq 2\pi$. On voit que le rayon de cette surface s’annule dans le plan d’équation $\vphi = \frac{\pi}{2}$, soit $x=0$. Un tel plan est appelé plan nodal, et est un plan d’antisymétrie de la fonction d’onde. Dans le plan d’équation $\theta = \frac{\pi}{2}$, soit $z=0$, on peut représenter plus simplement l’orbitale atomique dans un plan par la courbe d’équation $r = \cte  \times |\cos\vphi|$ en coordonnées polaires. On observe deux lobes dans la direction $(Ox)$.
  
  \end{exemple}

\end{document}
